% page 503 du fac-simile (image p-502 du scan)
% bloc 167.6 x 246.3 mm ; marge gauche 34.4 mm
\pgc{12.700mm}{0.000mm}{
\l{\cel{54}495.}
\l{}
\l{restar. (netrans.) Esar ankore prezenta pos ke onu deprenis un}
\l{\cel{3}o plura de lua parti. - DEFIS.}
\l{}
\l{restaurar. (trans.) (pri edifico, statuo, e c.) Igar en bona}
\l{\cel{3}stando itere. - DEFIRS.}
\l{}
\l{restitucar. (\sou{trans., ad.}) Retrodonar ad ulu (to quon onu prenis}
\l{\cel{3}de lu). - DEFIS.}
\l{}
\l{\sou{restorar}. (\sou{trans.)}\cel{2}Donar (ad ulu) la manjebli e drinkebli necesa}
\l{\cel{3}por ke lu ne plus hungrez nek durstez. - DEFIRS.}
\l{}
\l{\sou{restriktar}. (\sou{trans., ad)}. Retroduktar aden limiti plu streta.}
\l{\cel{3}- DEFIRS.}
\l{}
\l{\sou{reto}.\cel{2}I. Texuro ek mashi interligita per nodi plu o min inter-}
\l{\cel{3}distanta, facita ek kordeto, filo, e c., per bastono cilindra}
\l{\cel{3}e naveto. - II. Produkturo reto-forma ek filo, silko, filo ora}
\l{\cel{3}od arjenta. - III. Interplekturo de linei (voyi, kanali, rel-}
\l{\cel{3}voyi, e c.) (EFIS.}
\l{}
\l{\sou{retablo}. I. (\sou{tekn.}) Parto dopa di altaro (en kirko) qua supersa-}
\l{\cel{3}lias la tablo. - II. Panelo kontre qua, en kirko, apogesas al-}
\l{\cel{3}taro e quan ornas, en la kazi maxim multa, pikturo, skulturo,}
\l{\cel{3}en la parto qua superstacas la altaro. - FS.}
\l{}
\l{\sou{retenar}. (\sou{trans.}) I. Ne lasar de-irar (persono o kozo) ; impedar}
\l{\cel{3}(persono) agar impulsate da pasiono ; ne retrodonar (kozo),}
\l{\cel{3}ne restitucar (kozo). - II. Rezervigar por su (da ulu). - EFIS}
\l{}
\l{\sou{retiario.}\cel{2}(\sou{historio di la Roma antiqua)}. Gladiatoro di qua la}
\l{\cel{3}armo esis reto per qua lu probis enplektar sua adverso. - EFIS}
\l{}
\l{\sou{retikulo}. I. (\sou{optiko).}\cel{2}Fili interkrucumanta quin onu lokizas sur}
\l{\cel{3}la objektalo di la teleskopi, por determinar ula punti di la}
\l{\cel{3}feldo di ta aparati. - II. (\sou{biol.})\cel{1}\sou{Retikulo nukleala}\cel{1}:\cel{2}Parto}
\l{\cel{3}diferenciita di la nukleo celula, qua konsistas ye substanco}
\l{\cel{3}proteinala komplexa (\sou{linino)}, e qua rezultas de la pelotesko}
\l{\cel{3}di filamento unika e kontinua. -}
\l{}
\l{\sou{retino}. (\sou{anat}.)Membrano quan formacas, sur la fundo di la okulo,}
\l{\cel{3}la expanseso di la nervo vidala, sur qua venas facar su la}
\l{\cel{3}imajo de la objekti. - DEFIRS.}
\l{}
\l{\sou{retoro}. I. (\sou{epoki antiqua, en Roma e Grekia}).Profesoro pri elo-}
\l{\cel{3}quenteso. - II. (\sou{nun})\cel{2}Oratoro qua probas nur facar bela frazi.}
\l{\cel{3}- DEFIS.}
\l{}
\l{retorto. (\sou{kemio)}. Recipiento vitra, gresa, kun kolo streta e}
\l{\cel{3}kurvigita, uzata da la kemiisti. - II. (\sou{tekn}.)Tre granda reci-}
\l{\cel{3}piento ek tero refraktara, en qua onu varmigas min-karbono por}
\l{\cel{3}la fabriko di lum-gaso. - DEIRS.}
\l{}
\l{retraktar. (\sou{trans.})\cel{2}Retroprenar, nihiligar, explicite, formale}
\l{\cel{3}(to quon onu dicis, o plendo ad autoritato judiciala). - EFIS}
}
